\chapter{Metodología y datos}
\label{cap:metodologia}

Este capítulo describe el diseño experimental del trabajo: el dataset utilizado, los modelos
evaluados, las estrategias de prompting implementadas, el sistema de evaluación automática
y el protocolo de análisis cualitativo.

% ─────────────────────────────────────────────────────────────
\section{Dataset: Linguini}
\label{sec:dataset}
% ─────────────────────────────────────────────────────────────

Los experimentos utilizan puzzles extraídos del benchmark \textsc{Linguini}
\citep{sanchez2024linguini}, que provienen del corpus de la Olimpiada Internacional de
Lingüística (IOL). Se seleccionaron exclusivamente puzzles de tipo
\textit{sequence transduction} (traducción), que es la categoría más adecuada para evaluar
razonamiento metalingüístico inductivo y la única que permite comparación directa con los
trabajos de referencia \citep{sanchez2024linguini,zhu2025evaluating}.

\subsection{Criterios de selección}

Los criterios de selección de los puzzles fueron los siguientes:
(1)~pertenecer a familias lingüísticas distintas entre sí, para maximizar la diversidad
tipológica;
(2)~disponer de al menos ocho pares de entrenamiento, para que la estrategia incremental
tenga margen suficiente para construir conocimiento progresivo;
(3)~disponer de al menos tres pares de test, para evitar que las métricas dependan de una
única respuesta; y
(4)~no presentar indicios claros de conocimiento previo por parte de los modelos evaluados,
según el control exploratorio M0 (véase la Sección~\ref{subsec:m0}).

Se seleccionaron puzzles en ambas direcciones de traducción: de la lengua desconocida al
inglés (\textit{lang2en}) y del inglés a la lengua desconocida (\textit{en2lang}). Esta doble
dirección permite analizar si las estrategias de prompting son igualmente eficaces cuando el
modelo debe producir morfología de una lengua que no conoce, en lugar de interpretar
significado.

\subsection{Puzzles seleccionados}

El conjunto final comprende 16 puzzles en 11 lenguas y 10 familias lingüísticas distintas,
cubriendo regiones geográficas de Asia, Oceanía, América y África. La
Tabla~\ref{tab:puzzles} recoge las características de cada puzzle.

\begin{table}[h]
	\centering
	\caption{Puzzles utilizados en los experimentos. Dir.: dirección de traducción
		($\rightarrow$EN = lengua desconocida al inglés; EN$\rightarrow$ = inglés a lengua
		desconocida).}
	\label{tab:puzzles}
	\resizebox{\textwidth}{!}{%
		\begin{tabular}{llllccc}
			\toprule
			\textbf{ID} & \textbf{Lengua} & \textbf{Familia} & \textbf{Región} &
			\textbf{Dir.} & \textbf{Train} & \textbf{Test} \\
			\midrule
			\texttt{012018020100} & Hakhun         & Sino-Tibetana    & India          & $\rightarrow$EN  & 10 & 6 \\
			\texttt{012018020200} & Hakhun         & Sino-Tibetana    & India          & EN$\rightarrow$  & 10 & 4 \\
			\texttt{012022030100} & N|uuki         & Tuu              & Sudáfrica      & $\rightarrow$EN  & 16 & 6 \\
			\texttt{012022030200} & N|uuki         & Tuu              & Sudáfrica      & EN$\rightarrow$  & 16 & 6 \\
			\texttt{012019010100} & Yonggom        & Trans-New Guinea & Papua-NG       & $\rightarrow$EN  & 11 & 6 \\
			\texttt{012019010200} & Yonggom        & Trans-New Guinea & Papua-NG       & EN$\rightarrow$  & 11 & 4 \\
			\texttt{012015040100} & Wambaya        & Mirndi           & Australia      & $\rightarrow$EN  & 17 & 6 \\
			\texttt{012015040200} & Wambaya        & Mirndi           & Australia      & EN$\rightarrow$  & 17 & 5 \\
			\texttt{012023020100} & Apurinã        & Arawakan         & Brasil         & $\rightarrow$EN  & 16 & 3 \\
			\texttt{012023020300} & Apurinã        & Arawakan         & Brasil         & EN$\rightarrow$  & 16 & 6 \\
			\texttt{012023030100} & Coastal Marind & Trans-New Guinea & Indonesia      & $\rightarrow$EN  & 17 & 5 \\
			\texttt{012012010200} & Dyirbal        & Pama-Nyungan     & Australia      & $\rightarrow$EN  & 16 & 3 \\
			\texttt{012016030100} & Kunuz Nubian   & Nilo-Saharan     & Egipto/Sudán   & $\rightarrow$EN  & 10 & 5 \\
			\texttt{012008050100} & Inuktitut      & Eskimo-Aleut     & Canadá         & $\rightarrow$EN  & 12 & 6 \\
			\texttt{012005010100} & Tzotzil        & Maya             & México         & $\rightarrow$EN  & 15 & 4 \\
			\texttt{012006010100} & Lakhota        & Sioux            & EE.UU.         & $\rightarrow$EN  & 13 & 4 \\
			\bottomrule
		\end{tabular}%
	}
\end{table}

Cinco lenguas están representadas en ambas direcciones (Hakhun, N|uuki, Yonggom, Wambaya
y Apurinã), lo que permite comparar directamente el rendimiento de cada estrategia en
función de la dirección de traducción. Las lenguas restantes, con sus correspondientes
familias, permiten analizar si los patrones observados dependen de la tipología lingüística o
son generalizables.

\subsection{Formato de los datos}

Los puzzles de \textsc{Linguini} se distribuyen en formato JSONL. Para este trabajo se
implementó un conversor que transforma cada problema al formato JSON del sistema
experimental, con los campos \texttt{id}, \texttt{language}, \texttt{language\_family},
\texttt{translation\_direction}, \texttt{train\_pairs} y \texttt{test\_pairs}. Este formato
estandarizado permite añadir nuevos puzzles de forma sencilla y facilita la reproducibilidad
de los experimentos.

% ─────────────────────────────────────────────────────────────
\section{Modelos evaluados}
\label{sec:modelos}
% ─────────────────────────────────────────────────────────────

Se evaluaron tres modelos de lenguaje de gran escala de distintos tamaños, todos accesibles
mediante APIs gratuitas. Esta elección reduce el coste computacional y facilita la
reproducibilidad del trabajo. La Tabla~\ref{tab:modelos} recoge sus características
principales.

\begin{table}[h]
	\centering
	\caption{Modelos evaluados.}
	\label{tab:modelos}
	\resizebox{\textwidth}{!}{%
		\begin{tabular}{lllcll}
			\toprule
			\textbf{Denominación} & \textbf{Modelo} & \textbf{Organización} &
			\textbf{Parámetros} & \textbf{Tier} & \textbf{API} \\
			\midrule
			Llama-8B     & Llama 3.1 8B Instant     & Meta   & 8B   & Débil  & Groq \\
			Llama-70B    & Llama 3.3 70B Versatile  & Meta   & 70B  & Medio  & Groq \\
			GPT-OSS-120B & GPT-OSS 120B              & OpenAI & 120B & Fuerte & OpenRouter \\
			\bottomrule
		\end{tabular}%
	}
\end{table}

Los tres modelos son transformers \textit{decoder-only}, por lo que la diferencia principal
entre ellos es el tamaño (número de parámetros) y el proceso de preentrenamiento. La
elección de dos modelos de la familia Llama en tamaños diferentes (8B y 70B) permite
estudiar el efecto del tamaño manteniendo una arquitectura comparable. El tercer modelo
(GPT-OSS-120B) pertenece a una organización distinta y tiene un proceso de
preentrenamiento diferente, lo que añade diversidad a la comparación y permite analizar si
los patrones observados son específicos de los modelos Llama o generalizables.

Todos los modelos se emplean con temperatura cero (\texttt{temperature=0}) para reducir
la variabilidad entre ejecuciones, y sin ejemplos adicionales del formato de respuesta
esperado más allá de los propios del puzzle.

% ─────────────────────────────────────────────────────────────
\section{Estrategias de prompting}
\label{sec:estrategias}
% ─────────────────────────────────────────────────────────────

Se diseñó un control exploratorio previo, identificado como M0, y cinco estrategias de
traducción, identificadas como M1 a M5. Las cinco estrategias de traducción comparten una
instrucción metalingüística base. Además, se plantea una extensión futura de M2 basada en
la exploración exhaustiva de combinaciones de ejemplos.

\subsection{Instrucción metalingüística base}

Las estrategias M1 a M5 parten de una instrucción común que establece el marco de
razonamiento del modelo: actuar como analista lingüístico que nunca ha visto la lengua
objetivo, razonar exclusivamente a partir de los ejemplos proporcionados y tratar todas las
hipótesis como provisionales y revisables. Específicamente, se le indica que no utilice
conocimiento previo sobre la lengua o su familia, que construya explícitamente un
diccionario y reglas gramaticales, y que muestre su razonamiento antes de llegar a ninguna
conclusión. Este enfoque se inspira en el prompting metalingüístico de
\citet{zhu2025evaluating} y en los principios del \textit{chain-of-thought}
\citep{wei2022chain}.

\subsection{M0: Control exploratorio previo}
\label{subsec:m0}

Antes de ejecutar las estrategias de traducción, M0 intenta detectar posibles indicios de
conocimiento previo de la lengua del puzzle. Se presenta al modelo una única frase del
conjunto de test sin ningún contexto y se le pide que intente traducirla. Si el modelo produce
una traducción con chrF++ igual o superior a 20, se considera que podría existir conocimiento
previo y el experimento se marca para revisión manual.

Este control no permite demostrar que el modelo desconozca completamente la lengua. Su
objetivo es detectar casos potencialmente problemáticos antes de interpretar los resultados
de M1--M5.

\begin{figure}[ht]
	\centering
	\resizebox{\textwidth}{!}{%
		\begin{tikzpicture}[node distance=2cm, auto]
		\node (titulo) [text width=12cm, align=center]
		{\textbf{M0 --- Control exploratorio previo}\\
			\small ¿Existen indicios de conocimiento previo de la lengua?};
		
		\node (entrada) [rectverde, below of=titulo, yshift=-0.4cm, text width=4.4cm]
		{\textbf{Una frase de test}\\sin contexto de entrenamiento};
		
		\node (modelo) [rectvioleta, right of=entrada, xshift=4.8cm, text width=3cm]
		{\textbf{Modelo}\\traducción libre};
		
		\node (salida) [rectamarillo, right of=modelo, xshift=5cm, text width=5cm]
		{\textbf{chrF++ inferior a 20}\\
			no se detectan indicios claros\\
			$\rightarrow$ continuar con M1--M5};
		
		\node (nota) [rectnaranja, below of=modelo, yshift=-1.4cm, text width=9.5cm]
		{Si chrF++ es igual o superior a 20, el caso se marca para revisión manual.
			El control es heurístico.};
		
		\draw[arrow] (entrada) -- (modelo);
		\draw[arrow] (modelo) -- (salida);
		\end{tikzpicture}%
	}
	\caption{Funcionamiento del control exploratorio previo M0.}
	\label{fig:metodologia-m0}
\end{figure}

\subsection{M1: Baseline}
\label{subsec:baseline}

La estrategia baseline proporciona al modelo todos los pares de entrenamiento del puzzle
en un único prompt y le pide que genere directamente el diccionario, las reglas gramaticales
y las traducciones de las frases de test. Esta estrategia replica el prompting directo
utilizado como referencia en \citet{sanchez2024linguini} y sirve como punto de comparación
para las demás estrategias.

\begin{figure}[ht]
	\centering
	\resizebox{\textwidth}{!}{%
		\begin{tikzpicture}[node distance=2cm, auto]
		\node (titulo) [text width=12cm, align=center]
		{\textbf{M1 --- Baseline}\\
			\small Todos los ejemplos de entrenamiento se presentan en un único prompt};
		
		\node (entrada) [rectverde, below of=titulo, yshift=-0.4cm, text width=5cm]
		{\textbf{$n$ pares de entrenamiento}\\
			$f_1 \rightarrow t_1$\\
			$f_2 \rightarrow t_2$\\
			$\cdots$\\
			$f_n \rightarrow t_n$\\[1mm]
			+ instrucción metalingüística\\
			+ frases de test sin respuesta};
		
		\node (modelo) [rectvioleta, right of=entrada, xshift=5cm, text width=3cm]
		{\textbf{Modelo}\\una llamada a la API};
		
		\node (salida) [rectamarillo, right of=modelo, xshift=5cm, text width=5.2cm]
		{\textbf{Respuesta única}\\
			\texttt{DICTIONARY: ...}\\
			\texttt{GRAMMAR RULES: ...}\\
			\texttt{TRANSLATIONS: ...}};
		
		\draw[arrow] (entrada) -- (modelo);
		\draw[arrow] (modelo) -- (salida);
		\end{tikzpicture}%
	}
	\caption{Funcionamiento de la estrategia baseline M1.}
	\label{fig:metodologia-m1}
\end{figure}

\subsection{M2: Presentación incremental acumulativa}
\label{subsec:incremental}

La estrategia incremental constituye la aportación metodológica principal de este trabajo.
En lugar de presentar todos los ejemplos de golpe, se presentan prefijos de tamaño
creciente: primero las dos primeras frases, luego las tres primeras y así sucesivamente hasta
presentar las $n$ frases de entrenamiento, para un total de $\min(n-1, 15)$ pasos.

En cada paso, el modelo recibe el subconjunto de frases junto con su hipótesis acumulada
del paso anterior (diccionario y reglas) y debe realizar cuatro operaciones en orden:

\begin{enumerate}
	\item \textbf{Análisis} (\textsc{step~a}): razonar explícitamente sobre los nuevos
	ejemplos antes de actualizar ninguna hipótesis, identificando morfemas y detectando
	contradicciones con el conocimiento previo.
	
	\item \textbf{Diccionario revisado} (\textsc{step~b}): actualizar el diccionario añadiendo
	entradas nuevas, corrigiendo entradas imprecisas y eliminando entradas refutadas
	(marcadas como \texttt{[DELETED: razón]}).
	
	\item \textbf{Reglas revisadas} (\textsc{step~c}): actualizar el conjunto de reglas
	gramaticales con la misma lógica.
	
	\item \textbf{Traducciones} (\textsc{step~d}): traducir las frases de test usando el
	diccionario y las reglas actualizados.
\end{enumerate}

La posibilidad explícita de eliminar hipótesis erróneas es un elemento diferencial respecto
a otras estrategias incrementales, en las que el conocimiento acumulado tiende únicamente
a crecer. Se fundamenta en la observación de que los modelos pueden mantener hipótesis
incorrectas incluso cuando reciben evidencia contradictoria si el prompt no les obliga
explícitamente a revisarlas \citep{zhu2025evaluating}.

Para la comparación principal entre metodologías se utiliza el último paso de M2, en el que
el modelo ha recibido todos los ejemplos. El mejor paso observado según chrF++ se conserva
como información diagnóstica para analizar la evolución y la estabilidad del aprendizaje
incremental.

\begin{figure}[ht]
	\centering
	\resizebox{\textwidth}{!}{%
		\begin{tikzpicture}[node distance=1.5cm, auto]
		\node (titulo) [text width=13cm, align=center]
		{\textbf{M2 --- Presentación incremental acumulativa}\\
			\small En cada paso se añade una frase nueva y se revisa el conocimiento acumulado};
		
		\node (entrada1) [rectverde, below of=titulo, yshift=-0.5cm, text width=4.7cm]
		{\textbf{Paso 1}\\
			$f_1$, $f_2$};
		
		\node (modelo1) [rectvioleta, right of=entrada1, xshift=4.3cm, text width=3cm]
		{\textbf{Modelo}\\analiza e infiere};
		
		\node (salida1) [rectamarillo, right of=modelo1, xshift=5cm, text width=5.5cm]
		{\textbf{Diccionario v1 + reglas v1}\\
			traducciones de test};
		
		\node (entrada2) [rectverde, below of=entrada1, yshift=-1.1cm, text width=4.7cm]
		{\textbf{Paso 2}\\
			$f_1$, $f_2$, $f_3$\\
			+ diccionario v1 + reglas v1};
		
		\node (modelo2) [rectvioleta, right of=entrada2, xshift=4.3cm, text width=3cm]
		{\textbf{Modelo}\\analiza y revisa};
		
		\node (salida2) [rectamarillo, right of=modelo2, xshift=5cm, text width=5.5cm]
		{\textbf{Diccionario v2 + reglas v2}\\
			confirma, corrige o elimina hipótesis\\
			+ traducciones de test};
		
		\node (puntos) [below of=entrada2, yshift=-0.7cm]
		{\Large $\vdots$};
		
		\node (entradan) [rectverde, below of=puntos, yshift=-0.7cm, text width=4.7cm]
		{\textbf{Paso final}\\
			$f_1$, $f_2$, \ldots, $f_n$\\
			+ conocimiento acumulado};
		
		\node (modelon) [rectvioleta, right of=entradan, xshift=4.3cm, text width=3cm]
		{\textbf{Modelo}\\revisión final};
		
		\node (salidan) [rectamarillo, right of=modelon, xshift=5cm, text width=5.5cm]
		{\textbf{Diccionario final + reglas finales}\\
			traducciones de test del último paso};
		
		\draw[arrow] (entrada1) -- (modelo1);
		\draw[arrow] (modelo1) -- (salida1);
		\draw[arrow] (entrada2) -- (modelo2);
		\draw[arrow] (modelo2) -- (salida2);
		\draw[arrow] (entradan) -- (modelon);
		\draw[arrow] (modelon) -- (salidan);
		\end{tikzpicture}%
	}
	\caption{Funcionamiento de la estrategia incremental acumulativa M2.}
	\label{fig:metodologia-m2}
\end{figure}

\subsection{M2.2: Exploración combinatoria por tamaños crecientes}
\label{subsec:incremental-combinatoria}

Como ampliación de M2, se plantea una variante que recorre todas las combinaciones posibles
de frases de entrenamiento. Primero se evalúan todas las parejas, después todas las ternas y
así sucesivamente hasta llegar al conjunto completo. Esta variante permitirá estudiar si el
resultado depende del orden concreto en el que aparecen los ejemplos.

La implementación y evaluación completa de esta estrategia se incorporará en una fase
posterior del trabajo.

\begin{figure}[ht]
	\centering
	\resizebox{\textwidth}{!}{%
		\begin{tikzpicture}[node distance=1.1cm, auto]
		\node (titulo) [text width=14cm, align=center]
		{\textbf{M2.2 --- Exploración combinatoria por tamaños crecientes}\\
			\small Ejemplo con $n=4$ frases de entrenamiento};
		
		\node (k2) [rectverde, below of=titulo, yshift=-0.5cm, text width=4.8cm]
		{\textbf{$k=2$: $\binom{4}{2}=6$ combinaciones}\\
			$(f_1,f_2)$ \quad $(f_1,f_3)$ \quad $(f_1,f_4)$\\
			$(f_2,f_3)$ \quad $(f_2,f_4)$ \quad $(f_3,f_4)$};
		
		\node (k3) [rectverde, below of=k2, yshift=-0.7cm, text width=4.8cm]
		{\textbf{$k=3$: $\binom{4}{3}=4$ combinaciones}\\
			$(f_1,f_2,f_3)$ \quad $(f_1,f_2,f_4)$\\
			$(f_1,f_3,f_4)$ \quad $(f_2,f_3,f_4)$};
		
		\node (k4) [rectverde, below of=k3, yshift=-0.7cm, text width=4.8cm]
		{\textbf{$k=4$: $\binom{4}{4}=1$ combinación}\\
			$(f_1,f_2,f_3,f_4)$};
		
		\node (modelo) [rectvioleta, right of=k3, xshift=5.1cm, text width=3.2cm]
		{\textbf{Modelo}\\
			una llamada por combinación};
		
		\node (salida) [rectamarillo, right of=modelo, xshift=5cm, text width=5.2cm]
		{\textbf{Resultado por iteración}\\
			diccionario + reglas\\
			+ traducciones de test};
		
		\node (coste) [rectnaranja, below of=salida, yshift=-1cm, text width=6.3cm]
		{\textbf{Coste computacional}\\
			Con $n=10$: $\sum_{k=2}^{10}\binom{10}{k}=1013$ iteraciones};
		
		\draw[arrow] (k2.east) -- (modelo.north west);
		\draw[arrow] (k3) -- (modelo);
		\draw[arrow] (k4.east) -- (modelo.south west);
		\draw[arrow] (modelo) -- (salida);
		\end{tikzpicture}%
	}
	\caption{Extensión planificada de M2 mediante la exploración exhaustiva de
		combinaciones de ejemplos.}
	\label{fig:metodologia-m22}
\end{figure}

\subsection{M3: Step-by-step lingüístico}
\label{subsec:stepbystep}

La estrategia step-by-step está inspirada en \citet{zhu2025evaluating} y organiza la
presentación de los ejemplos según una jerarquía lingüística de cuatro etapas:

\begin{enumerate}
	\item \textbf{Léxico} (35\,\% de los ejemplos): frases simples para identificar
	vocabulario básico y orden de palabras fundamental.
	\item \textbf{Fonología} (15\,\% de los ejemplos): frases que exhiben variaciones
	fonológicas (alternancias vocálicas, alófonos, asimilaciones).
	\item \textbf{Morfosintaxis} (30\,\% de los ejemplos): frases con concordancia de
	persona, número, tiempo y aspecto.
	\item \textbf{Sintaxis} (20\,\% de los ejemplos): frases con estructuras complejas
	(negación, interrogación, subordinación).
\end{enumerate}

Esta partición es heurística cuando los puzzles carecen de anotación lingüística explícita,
como ocurre en \textsc{Linguini}. La asignación de los ejemplos a cada etapa se realiza según
su posición en el conjunto de entrenamiento y no garantiza que cada bloque contenga
exclusivamente el fenómeno indicado. Al igual que en M2, el modelo acumula diccionario y
reglas entre etapas y puede corregir hipótesis previas.

\begin{figure}[ht]
	\centering
	\resizebox{\textwidth}{!}{%
		\begin{tikzpicture}[node distance=2cm, auto]
		\node (titulo) [text width=13cm, align=center]
		{\textbf{M3 --- Step-by-step lingüístico}\\
			\small Cuatro etapas en orden fijo con acumulación de conocimiento};
		
		\node (lexico) [rectverde, below of=titulo, yshift=-0.5cm, text width=3cm]
		{\textbf{1. Léxico}\\35\,\%};
		
		\node (fonologia) [rectvioleta, right of=lexico, xshift=3cm, text width=3cm]
		{\textbf{2. Fonología}\\15\,\%};
		
		\node (morfosintaxis) [rectamarillo, right of=fonologia, xshift=3cm, text width=3cm]
		{\textbf{3. Morfosintaxis}\\30\,\%};
		
		\node (sintaxis) [rectnaranja, right of=morfosintaxis, xshift=3cm, text width=3cm]
		{\textbf{4. Sintaxis}\\20\,\%};
		
		\node (salida) [rectamarillo, below of=fonologia, xshift=3cm, yshift=-1.2cm,
		text width=9.5cm]
		{\textbf{En cada etapa, el modelo recibe:}\\
			ejemplos de la etapa actual + diccionario y reglas acumulados\\
			$\rightarrow$ genera conocimiento revisado y traducciones de test};
		
		\node (nota) [rectnaranja, below of=salida, yshift=-1cm, text width=9.5cm]
		{\textbf{Limitación:} la partición es heurística y no está curada manualmente
			según el nivel lingüístico real.};
		
		\draw[arrow] (lexico) -- (fonologia);
		\draw[arrow] (fonologia) -- (morfosintaxis);
		\draw[arrow] (morfosintaxis) -- (sintaxis);
		\draw[arrow] (morfosintaxis) -- (salida);
		\end{tikzpicture}%
	}
	\caption{Funcionamiento de la estrategia step-by-step lingüística M3.}
	\label{fig:metodologia-m3}
\end{figure}

\subsection{M4: Inspiración humana}
\label{subsec:humana}

La estrategia de inspiración humana proporciona al modelo el razonamiento explícito de una
persona resolviendo un puzzle diferente, para guiar el proceso inductivo. El razonamiento
utilizado como ejemplo es la resolución detallada del puzzle de Hakhun (véase la
Sección~\ref{sec:cualitativo}), que incluye los patrones identificados en cada frase, las
hipótesis provisionales formuladas y cómo se corrigen o confirman.

El modelo recibe este razonamiento como referencia de proceso, no de contenido: se le indica
explícitamente que no transfiera vocabulario ni reglas del ejemplo. El objetivo es que el
modelo imite la estrategia de análisis (identificar invariantes, aislar morfemas y revisar
hipótesis progresivamente) aplicándola a una nueva lengua.

Para evitar una fuga directa de información, M4 no se ejecuta sobre ningún puzzle de
Hakhun, independientemente de la dirección de traducción.

\begin{figure}[ht]
	\centering
	\resizebox{\textwidth}{!}{%
		\begin{tikzpicture}[node distance=2cm, auto]
		\node (titulo) [text width=13cm, align=center]
		{\textbf{M4 --- Inspiración humana}\\
			\small Una resolución humana guía el proceso inductivo aplicado a otra lengua};
		
		\node (humano) [rectamarillo, below of=titulo, yshift=-0.5cm, text width=4.7cm]
		{\textbf{Resolución humana de Hakhun}\\
			análisis paso a paso\\
			identificación de patrones\\
			revisión de hipótesis};
		
		\node (nuevo) [rectverde, below of=humano, yshift=-0.8cm, text width=4.7cm]
		{\textbf{Puzzle nuevo}\\
			pares de entrenamiento\\
			+ frases de test};
		
		\node (modelo) [rectvioleta, right of=nuevo, xshift=5cm, text width=3.4cm]
		{\textbf{Modelo}\\
			imita el proceso,\\
			no el contenido};
		
		\node (salida) [rectamarillo, right of=modelo, xshift=5cm, text width=5.3cm]
		{\textbf{Respuesta final}\\
			diccionario + reglas\\
			+ traducciones de test};
		
		\node (nota) [rectnaranja, below of=modelo, yshift=-1.5cm, text width=8.8cm]
		{M4 no se ejecuta sobre ningún puzzle de Hakhun para evitar fuga de
			información.};
		
		\draw[arrow] (humano.east) -- (modelo.north west);
		\draw[arrow] (nuevo) -- (modelo);
		\draw[arrow] (modelo) -- (salida);
		\end{tikzpicture}%
	}
	\caption{Funcionamiento de la estrategia de inspiración humana M4.}
	\label{fig:metodologia-m4}
\end{figure}

\subsection{M5: Autocorrección iterativa}
\label{subsec:autocorreccion}

La estrategia de autocorrección iterativa está inspirada en el paradigma
\textit{Self-Refine} \citep{madaan2023selfrefine} y opera en dos turnos. En el primer turno,
el modelo realiza una traducción completa siguiendo el mismo prompt que M1. En el segundo
turno, recibe sus propias traducciones y el diccionario y reglas que infirió, y se le pide que
identifique contradicciones internas explícitas antes de producir una versión corregida. La
instrucción es deliberadamente precisa: debe citar la regla o entrada de diccionario que cada
traducción viola o que no está cubierta, y solo entonces producir las traducciones finales.

Esta estrategia evalúa si el razonamiento metalingüístico del propio modelo es suficiente
para detectar y corregir errores de coherencia sin recibir información adicional sobre la
lengua.

\begin{figure}[ht]
	\centering
	\resizebox{\textwidth}{!}{%
		\begin{tikzpicture}[node distance=1.7cm, auto]
		\node (titulo) [text width=13cm, align=center]
		{\textbf{M5 --- Autocorrección iterativa}\\
			\small Dos turnos: el modelo revisa su propio borrador buscando contradicciones};
		
		\node (entrada1) [rectverde, below of=titulo, yshift=-0.5cm, text width=4.7cm]
		{\textbf{Turno 1}\\
			pares de entrenamiento\\
			+ frases de test\\
			+ instrucción metalingüística};
		
		\node (modelo1) [rectvioleta, right of=entrada1, xshift=4.5cm, text width=3cm]
		{\textbf{Modelo}\\turno 1};
		
		\node (borrador) [rectamarillo, right of=modelo1, xshift=5cm, text width=5cm]
		{\textbf{Borrador inicial}\\
			diccionario + reglas\\
			+ traducciones};
		
		\node (entrada2) [rectverde, below of=entrada1, yshift=-1.1cm, text width=4.7cm]
		{\textbf{Turno 2}\\
			mismos ejemplos\\
			+ borrador inicial\\
			+ solicitud de crítica};
		
		\node (modelo2) [rectvioleta, right of=entrada2, xshift=4.5cm, text width=3cm]
		{\textbf{Modelo}\\turno 2};
		
		\node (final) [rectamarillo, right of=modelo2, xshift=5cm, text width=5cm]
		{\textbf{Respuesta revisada}\\
			crítica explícita\\
			+ diccionario y reglas revisados\\
			+ traducciones finales};
		
		\draw[arrow] (entrada1) -- (modelo1);
		\draw[arrow] (modelo1) -- (borrador);
		\draw[arrow] (borrador.south) |- (entrada2.east);
		\draw[arrow] (entrada2) -- (modelo2);
		\draw[arrow] (modelo2) -- (final);
		\end{tikzpicture}%
	}
	\caption{Funcionamiento de la estrategia de autocorrección iterativa M5.}
	\label{fig:metodologia-m5}
\end{figure}

\FloatBarrier

% ─────────────────────────────────────────────────────────────
\section{Sistema experimental}
\label{sec:sistema}
% ─────────────────────────────────────────────────────────────

Se implementó un sistema experimental en Python (\texttt{experiment\_runner.py}) que
automatiza la ejecución de los experimentos, el cálculo de métricas y el registro de
resultados. Dado un puzzle y un modelo, el sistema ejecuta el control M0 y las estrategias
M1--M5, calcula las métricas en cada paso y guarda un log completo en formato JSON con las
hipótesis extraídas, las métricas calculadas, el diccionario y las reglas acumulados en cada
paso, y opcionalmente el prompt y la respuesta completa del modelo.

\subsection{Extracción de traducciones}

Un aspecto técnico relevante es la extracción de las traducciones de la respuesta libre del
modelo. Los modelos no siempre siguen el formato solicitado: algunos utilizan comillas
tipográficas (\texttt{\textquotedblleft\textquotedblright}), cursivas en Markdown
(\texttt{*texto*}), flechas unicode (\texttt{→}) u otros formatos. Se implementó un
extractor multicapa con expresiones regulares que cubre los patrones de salida observados
empíricamente, con filtros explícitos para descartar entradas de diccionario, instrucciones
de revisión y meta-análisis lingüístico que los modelos a veces incluyen en la sección de
traducciones. Este extractor fue validado de forma independiente comparando las métricas
calculadas contra un recálculo manual sobre los mismos datos.

\subsection{Parámetros de ejecución}

Las llamadas a los modelos se realizan con temperatura cero y un límite de 4096 tokens de
salida. Para respetar los límites de tasa de las APIs gratuitas, se introduce un retardo de
10 segundos entre los pasos consecutivos de M2 y de 20 segundos entre los dos turnos de M5.
El tiempo de espera máximo por llamada es de 150 segundos, configurado especialmente para
OpenRouter, cuya latencia con el modelo de 120B parámetros es significativamente mayor.

% ─────────────────────────────────────────────────────────────
\section{Métricas de evaluación}
\label{sec:metricas}
% ─────────────────────────────────────────────────────────────

Se utilizan cuatro métricas automáticas de evaluación, calculadas con la biblioteca
\texttt{sacrebleu} \citep{post2018sacrebleu}:

\textbf{BLEU} \citep{papineni2002bleu} calcula la precisión de n-gramas compartidos entre
la hipótesis del modelo y la referencia, con una penalización por brevedad. Su principal
limitación en este contexto es que requiere coincidencia exacta de palabras, lo que la hace
excesivamente estricta para lenguas con morfología compleja. Se incluye por compatibilidad
con la literatura existente ($\uparrow$ mejor, rango 0--100).

\textbf{chrF++} \citep{popovic2017chrf} es la métrica principal de este trabajo. Opera a
nivel de n-gramas de caracteres y de palabras (\texttt{word\_order=2}), lo que la hace más
tolerante con variaciones morfológicas y más adecuada para lenguas aglutinantes. Es la
métrica secundaria del propio benchmark \textsc{Linguini} \citep{sanchez2024linguini} y
la utilizada para analizar la evolución de las estrategias iterativas
($\uparrow$ mejor, rango 0--100).

\textbf{TER} \citep{snover2006ter} (\textit{Translation Edit Rate}) mide el número mínimo
de ediciones por palabra necesarias para transformar la hipótesis en la referencia. A
diferencia de BLEU y chrF++, es una métrica de error: un valor de 0 indica traducción
perfecta, y valores superiores a 100 indican que se necesitan más ediciones que palabras hay
en la referencia ($\downarrow$ mejor).

\textbf{Exact Match} (EM) cuenta el porcentaje de traducciones exactamente correctas,
ignorando mayúsculas y puntuación. Con pocos pares de test por puzzle (entre 3 y 6), cada
acierto representa entre un 17 y un 33\,\% del score, lo que hace esta métrica discreta pero
muy directamente interpretable ($\uparrow$ mejor, rango 0--100\,\%).

% ─────────────────────────────────────────────────────────────
\section{Análisis cualitativo}
\label{sec:cualitativo}
% ─────────────────────────────────────────────────────────────

Además de la evaluación automática, se realizó un análisis cualitativo comparando el proceso
de resolución humana con el del modelo en el puzzle de Hakhun (Sino-Tibetana, India;
10 frases de entrenamiento, 6 de test). El objetivo es identificar en qué aspectos del
razonamiento metalingüístico los modelos divergen de los humanos, más allá de lo que
capturan las métricas automáticas.

\subsection{Protocolo de resolución humana}

El protocolo consistió en resolver el puzzle anotando explícitamente cada paso del
razonamiento: qué patrones se identificaron primero, qué hipótesis se formularon, cuándo y
por qué se corrigieron, y cómo se aplicaron a las frases de test.

\subsection{Proceso de resolución observado}

El proceso humano siguió el siguiente orden: primero se identificaron los elementos
invariantes (\texttt{ne} como marcador de pregunta, presente en todas las frases), luego los
pronombres de sujeto (\texttt{ŋa}=I, \texttt{nɤ}=you(sg), \texttt{nirum}=we,
\texttt{tarum}=they), después los verbos (\texttt{lapkʰi}=see, \texttt{cʰam}=know,
\texttt{lan}=beat, \texttt{ka}=go, \texttt{ʒip}=sleep), y finalmente el sistema aspectual:
las terminaciones con \texttt{ʔ} o \texttt{tʰ-} marcan pasado, mientras que las demás
marcan presente. El marcador \texttt{kəmə} se identificó como marcador de transitividad sin
traducción directa al inglés. Con este protocolo se obtuvieron las seis traducciones de test
correctas.

\subsection{Comparación con el modelo}

La comparación con la respuesta del modelo en M1 revela un patrón consistente: el modelo
identifica correctamente el léxico básico (verbos y pronombres de sujeto), pero falla
sistemáticamente en el sistema aspectual (pasado frente a presente). El modelo menciona en
su razonamiento la posibilidad de que \texttt{ʔ} marque pasado, pero no lo aplica de forma
consistente a todas las formas verbales. Este resultado está en línea con el hallazgo de
\citet{zhu2025evaluating} de que los modelos rinden peor en puzzles con mayor complejidad
morfológica, ya que la concordancia de tiempo, aspecto y objeto en Hakhun es precisamente
el elemento más complejo del puzzle.

En M2, el modelo con subconjuntos pequeños produce traducciones degradadas por
insuficiencia de evidencia para deducir el sistema pronominal de objeto. El rendimiento
mejora en pasos intermedios, pero puede decaer en los pasos finales en los modelos más
débiles, posiblemente por sobrecarga del contexto del prompt. El humano mostró el patrón
opuesto: con solo dos frases ya identificó los elementos invariantes, y cada par adicional
refinó hipótesis previas sin degradar las ya establecidas.

Este contraste sugiere que la ventaja de los modelos grandes en tareas de traducción
metalingüística podría estar relacionada con una mayor capacidad para mantener coherencia
en contextos largos y con una base de conocimiento tipológico más amplia sobre la que
apoyar las inferencias.